%%%%%%%%%%  scan idx 669  |  folio 656  %%%%%%%%%%
\SGAPrintPage{656}%
From $s\ge1$ follows $G'\ne0$; consequently, $F'$ is not contained in $F(-1)$; that is, the section $\sigma$ is not contained in the image of the injection $H^0(F(-1))\to H^0(F)$, whence (i).

Assertion (ii) follows immediately from (i), taking into account that $h^0(F)=0$ clearly implies $h^0(F(-1))=0$, since $h^0(F(-1))\le h^0(F)$. In (iii), we find $0\le h^0(F(n))-h^0(F(n-1))\le h^0(G(n))$. Thus, for $n\ge n_0$, $h^0(F(-n_0))=h^0(F(-n))$ since $h^0(G(-n_0))=0$; hence $h^0(F(-n_0))=0$ by (ii). Consequently, $h^0(F(p-n_0))\le h^0(G(1-n_0))+\ldots+h^0(G(p-n_0))\le ph^0(G(p-n_0))$ for $p\ge1$; whence (iii).

\sgaStableTarget{sga6.e13.prop.6-6}\textbf{Proposition 6.6.} \emph{Under the conditions of 6.1, suppose that $F$ is a $(b_0,\ldots,b_r)$-sheaf. Then $F_q$ is a $\bigl((b_{q-1})^q,\ldots,(b_{q-1})^2,b_{q-1},b_q,\ldots,b_r\bigr)$-sheaf.}

Indeed, let $\sigma$ be a ``general'' section of ${O}_P(1)$; then $\sigma$ defines a commutative diagram with exact rows and columns
\[
\begin{tikzcd}[column sep=large,row sep=large]
&0 \arrow[d]&0 \arrow[d]&0 \arrow[d]&\\
0 \arrow[r]&N_q(-1) \arrow[r] \arrow[d]&N_q \arrow[r] \arrow[d]&G' \arrow[r] \arrow[d]&0\\
0 \arrow[r]&F(-1) \arrow[r] \arrow[d]&F \arrow[r] \arrow[d]&G \arrow[r] \arrow[d]&0\\
0 \arrow[r]&F_q(-1) \arrow[r] \arrow[d]&F_q \arrow[r] \arrow[d]&G'' \arrow[r] \arrow[d]&0\\
&0&0&0&
\end{tikzcd}
\]
where $G$ is a $(b_1,\ldots,b_r)$-sheaf (1.6 (i)) and where $G''=G_{q-1}$ (a question of depth).

%%%%%%%%%%  scan idx 670  |  folio 657  %%%%%%%%%%
\SGAPrintPage{657}%
Suppose that $q=1$. Then $\operatorname{dim{.}supp}(N_q)=0$; thus $H^1(N_q(-1))=0$ and $h^0(F_q(-1))\le h^0(F(-1))\le b_0$. Moreover, $G'=0$; hence $G\simeq G''$. Thus $F_1$ is indeed a $(b_0,\ldots,b_r)$-sheaf. In addition, 6.5 (ii) shows that $H^0(F_1(-b_0))=0$. If $q\ge2$, we may suppose that $G_{q-1}$ is a $((b_{q-1})^{q-1},\ldots,b_{q-1},\ldots,b_r)$-sheaf and that $H^0(G_{q-1}(-b_{q-1}))=0$, by induction on $q$. Then 6.5 shows that $H^0(F_q(-b_{q-1}))=0$ and $h^0(F_q(-1))\le (b_{q-1})^q$; whence the assertion.

\sgaStableTarget{sga6.e13.thm.6-7}\textbf{Theorem 6.7.} \emph{There exist two sequences $\{A_i(X_0,\ldots,X_i;Y)\}$ and $\{A_i^{(q)}(X_0,\ldots,X_q;Y)\}$ of polynomials with the following properties. Assume that the conditions of 6.1 are satisfied and that $F$ is a quotient of ${O}_P(-m)^{\oplus M}$ with $m\ge0$.}
\begin{enumerate}[label=(\roman*)]
\item \emph{If $F$ is a $(b_0,\ldots,b_r)$-sheaf, then, for $i=0,\ldots,r$,}
\[
   |a_i|\le A_{r-i}(b_i,\ldots,b_r;m).
\]
\item \emph{If $\chi(F_q(n))=\sum_{i=0}^{r} a_i^{(q)}\binom{n+i}{i}$ is the Hilbert polynomial of $F_q$, then, for $i=0,\ldots,q-1$,}
\[
   |a_i^{(q)}|\le A_{r-i}^{(r-q+1)}(c_{q-1},\ldots,c_r;m).
\]
\end{enumerate}
\emph{(Notice that $a_{q-1}^{(q)}\le a_{q-1}$ and $a_q^{(q)}=a_q,\ldots,a_r^{(q)}=a_r$ since $\operatorname{dim{.}supp}(N_q)\le q-1$).}

Indeed, let an exact sequence $0\to F(-1)\to F\to G\to0$ be such that $G$ is a $(b_1,\ldots,b_r)$-sheaf; arguing by induction on $r$, we may suppose that $A_0,\ldots,A_{r-1}$ have been defined, taking 1.7 into account. Then 6.4 (ii) makes the definition of $A_r$ clear, taking 1.8 into account. Finally, (ii) follows formally from 6.4 (iii), 6.6, and assertion (i) (where we take $F_q$ for $F$).

%%%%%%%%%%  scan idx 671  |  folio 658  %%%%%%%%%%
\SGAPrintPage{658}%
\sgaStableTarget{sga6.e13.cor.6-8}\textbf{Corollary 6.8. (Grothendieck).} \emph{Let $X$ be projective over noetherian $S$, let ${O}_X(1)$ be very ample for $X/S$, and let $\mathcal J$ be a family of classes of coherent sheaves on the fibers of $X/S$. Assume}:
\begin{enumerate}[label=(\alph*)]
\item \emph{There exists a coherent sheaf $E$ on $X$ such that $\mathcal J$ is contained in the family of classes of quotients of sheaves of the form $E_K$.}
\item[$(b_q)$] \emph{In the Hilbert polynomials $\chi(F_K(n))$ of the $F_K\in\mathcal J$, the coefficients in degrees $\ge q-1$ remain bounded above.}
\end{enumerate}
\emph{Then the $(F_K)_q$ (as $F_K$ ranges over $\mathcal J$) form a bounded family; moreover, the coefficients in $\chi(F_K(n))$ in degrees $\ge q-2$ remain bounded below.}

Indeed, we may clearly suppose that $X=\mathbb P^N_S$ and $E={O}_X(-m)^{\oplus M}$. Then the first assertion follows formally from 6.7 (ii) and 1.13; the second follows from 6.4 (ii), by induction using an exact sequence $0\to F(-1)\to F\to G\to0$.

\sgaStableTarget{sga6.e13.def.6-9}\textbf{Definition 6.9.} \emph{Let $k$ be a field. A \emph{special positive $k$-cycle of dimension $r$} is a projective $k$-scheme $X$, equipped with a very ample sheaf ${O}_X(1)$, which is a union of closed subschemes $X_j$ of dimension $r$, with $X_j$ obtained by a base change $\operatorname{Spec}(k)\to\operatorname{Spec}(k_j)$ as the inverse image of an \emph{integral} $k_j$-scheme $X'_j$ (equipped with ${O}_{X'_j}(1)$). The \emph{degree} of $X$ is the coefficient $a_r$ in the Hilbert polynomial $\chi({O}_X(n))=\sum a_i\binom{n+i}{i}$.}

\sgaStableTarget{sga6.e13.lem.6-10}\textbf{Lemma 6.10.} \emph{Let $k$ be an algebraically closed field and $X$ a special positive $k$-cycle of dimension $r$ and degree $d$. Then ${O}_X$ is a $(0,\ldots,0,d)$-sheaf.}

%%%%%%%%%%  scan idx 672  |  folio 659  %%%%%%%%%%
\SGAPrintPage{659}%
Indeed, let $X$ lie in $\mathbb P^N_k$. After replacing $k$ by the algebraic closure of a purely transcendental extension, which is allowed by 1.6 (i), we may suppose that there are $r(N+1)$ elements in $k$ which are algebraically independent over all the $k_j$. Then the $r$ successive ``generic'' hyperplane sections of $X$, defined by means of these elements, are still special positive $k$-cycles of degree $d$.

Arguing by induction on $r$, it therefore remains to note that, when $r=0$, we have $h^0({O}_X(-1))=d$, and, when $r\ge1$, we have $h^0({O}_X(-1))=0$. For this last point, we may suppose that $X$ itself comes from an integral scheme $X'$, by the injection ${O}_X\to \coprod{O}_{X_j}$; but then $h^0({O}_X(-1))=h^0({O}_{X'}(-1))$, which is $0$ since $H^0({O}_{X'})$ is a field.

\sgaStableTarget{sga6.e13.cor.6-11}\textbf{Corollary 6.11.} (i) \emph{Let $r,d$ be two integers, let $c_i=A_{r-i}(0,\ldots,0,d;0)$ for $i=0,\ldots,r$ $(6.7\ \mathrm{(i)})$, and let $p=P_r(c_0,\ldots,c_r)$ $(1.9)$. Let $k$ be an algebraically closed field, and $X$ a special positive $k$-cycle of dimension $\le r$ and degree $\le d$, with Hilbert polynomial $\chi({O}_X(n))=\sum_{i=0}^{r}a_i\binom{n+i}{i}$. Then, for $i=0,\ldots,r$, we have}
\[
   |a_i|\le c_i.
\]
\emph{Moreover, $X$ embeds as a subscheme of $\mathbb P^N_k$ with $N=d(r+1)-1$, defined by (at most) $\binom{N+p}{N}$ equations of degree $p$.}

(ii) (Chow) \emph{Let $S$ be a noetherian scheme. As $K$ ranges over the algebraically closed extensions of the $k(s)$, $s\in S$, the special positive $K$-cycles of bounded dimension and degree form a ``bounded family'' (either with the $X$ regarded as subschemes of a fixed $\mathbb P^N_S$, or in an evident abstract sense).}

%%%%%%%%%%  scan idx 673  |  folio 660  %%%%%%%%%%
\SGAPrintPage{660}%
Indeed, in (i), the first assertion follows from 6.7 (i); the second, from 6.10 and 1.8 (i), and from 6.2 (applied to $0\to I\to{O}_P\to{O}_X\to0$, $P=\mathbb P^N_k$) and 1.3. Finally, (ii) follows from (i) and 1.13.

\sgaStableTarget{sga6.e13.rem.6-12}\textbf{Remark 6.12.} Let $k$ be an algebraically closed field, and $X$ a projective $k$-scheme of pure dimension $r$, with no embedded associated prime cycle. If we allow arbitrary nilpotent elements in ${O}_X$, we can no longer bound the coefficients in the Hilbert polynomial $\chi({O}_X(n))=\sum_{i=0}^{r}a_i\binom{n+i}{i}$ in terms of $a_r$ alone.

For example, let $Z$ be a smooth curve of degree $d$ in $\mathbb P^3_k$. For every $n\gg0$, there exists a smooth surface $Y$ of degree $n$ which contains $Z$. Then $Z$ is a Cartier divisor on $Y$; put ${O}_{X_n}={O}_Y/{O}_Y(-2Z)$. Finally, we have the exact sequence $0\to{O}_Z(-Z)\to{O}_{X_n}\to{O}_Z\to0$ and the formula $\langle Z^2\rangle_Y=2p_a(Z)-2-(n-4)d$; hence $a_0=\chi({O}_{X_n}(-1))\to\infty$ ($n\to\infty$), although $a_1=2d$.

\sgaStableTarget{sga6.e13.cor.6-13}\textbf{Corollary 6.13.} \emph{Under the conditions of 6.1, consider an exact sequence}
\[
   {O}_P(-m_1)^{\oplus M_1}\longrightarrow
   {O}_P(-m)^{\oplus M}\xrightarrow{\alpha}F\longrightarrow0
\]
\emph{with $m_1\ge m\ge0$. Then, for $i=0,\ldots,N$,}
\[
   |a_i|\le A_{N-i}(0,\ldots,0,M;m_1)+M\binom{m}{N-i}.
\]

Indeed, put $I=\operatorname{Ker}(\alpha)$; clearly, $I$ is a $(0,\ldots,0,M)$-sheaf, and $\chi(I(n))=\sum_{i=0}^{N}\left[M\binom{-m-1+N-i}{N-i}-a_i\right]\binom{n+i}{i}$; whence the assertion by 6.7 (i), and the formula $\binom{-m-1+Q}{Q}=(-1)^Q\binom{m}{Q}$.

